\scanpage{48}
\begin{center}4 - 41 -\hfill n\textsuperscript{o} 392\end{center}

principal semi-local, dont les idéaux maximaux sont les $\underline{p}S^{-1}A$, où $\underline{p}\in Y$, les anneaux locaux correspondants étant les $A_{\underline{p}}$. On a donc $S^{-1}A/\underline{p}S^{-1}A=k(\underline{p})$, d'où $S^{-1}M/\underline{p}S^{-1}M=(M/\underline{p}M)\otimes_A S^{-1}A=M\otimes_A(A/\underline{p}\otimes_A S^{-1}A)=M\otimes_A k(\underline{p})=M(\underline{p})$ pour tout $\underline{p}\in Y$. En vertu du chap. II, § 1, n\textsuperscript{o} 2, prop. 6, il existe un élément $z/s\in S^{-1}M$, $(z\in M,s\in S)$, dont les images canoniques dans les $M(\underline{p})$ pour $\underline{p}\in Y$ sont toutes $\ne0$. Par définition de $S$, on a donc $z(\underline{p})\ne0$ pour tout $\underline{p}\in Y$. On peut en outre supposer que $y$ et $z$ sont \uline{linéairement indépendants} sur $K$. En effet, dans le cas contraire, considérons un élément $t\in M$ linéairement indépendant de $y$, (il en existe puisque $n\ge2$) ; prenons d'autre part un élément $a\ne0$ appartenant à $\bigcap_{\underline{p}\in Y}\underline{p}$, (qui n'est pas réduit à $0$ puisque $A$ est intègre), et posons $z'=z+at$ : il est clair que $y$ et $z'$ sont linéairement indépendants sur $K$ et que l'on a $z'(\underline{p})=z(\underline{p})\ne0$ pour tout $\underline{p}\in Y$.

Supposant donc $y$ et $z$ linéairement indépendants sur $K$, soit $Z$ l'ensemble des $\underline{p}\in P-Y$ tels que $y(\underline{p})$ et $z(\underline{p})$ soient linéairement dépendants sur $k(\underline{p})$ ; il résulte du lemme 7 que cet ensemble est \uline{fini}. Pour tout $\underline{p}\in Z$, on peut donc écrire $z(\underline{p})=\lambda(\underline{p})y(\underline{p})$ avec $\lambda(\underline{p})\in k(\underline{p})$. Or, on a $\operatorname{Card}(A/\underline{p})\ge2$ pour tout $\underline{p}\in P$ ; il résulte donc du lemme 6 qu'il existe $b\in A$ tel que, pour tout $\underline{p}\in Z$, l'image canonique de $b$ dans $A/\underline{p}$ soit distincte de $\lambda(\underline{p})$. Montrons alors que l'élément $x=z-by$ répond à la question ; il suffit, (en vertu du lemme 7 appliqué pour $m=1$), de vérifier que $x(\underline{p})\ne0$ pour \uline{tout} $\underline{p}\in P$. Or :

- si $\underline{p}\in Y$, on a $x(\underline{p})=z(\underline{p})\ne0$ par construction ;

- si $\underline{p}\in Z$, on a $x(\underline{p})=\mu_{\underline{p}}.y(\underline{p})$ avec $\mu_{\underline{p}}\ne0$ en vertu du choix de $b$, donc $x(\underline{p})\ne0$ puisque $y(\underline{p})\ne0$ ;

- si $\underline{p}\in P-(Y\cup Z)$, $y(\underline{p})$ et $z(\underline{p})$ sont linéairement indépendants, donc $x(\underline{p})\ne0$.

Ces lemmes étant établis, passons à la démonstration du th. 6. Raisonnons par récurrence sur $n$, le cas $n\le1$ étant trivial puisque $M$ lui-même est alors iso-

\margin{47}
